\documentclass[UTF8,12pt,a4paper,fontset=fandol]{ctexart}

\usepackage[a4paper,top=24mm,bottom=24mm,left=26mm,right=26mm]{geometry}
\usepackage{amsmath,amssymb,bm}
\usepackage{booktabs,longtable,array,enumitem}
\usepackage{xcolor,hyperref,fancyhdr,microtype}

\hypersetup{
  colorlinks=true,
  linkcolor=blue!55!black,
  urlcolor=blue!55!black
}

\setlength{\parindent}{2em}
\setlength{\parskip}{0.30em}
\setlist[itemize]{leftmargin=2.2em,itemsep=0.18em,topsep=0.30em}
\setlist[enumerate]{leftmargin=2.4em,itemsep=0.20em,topsep=0.30em}
\allowdisplaybreaks[2]

\pagestyle{fancy}
\setlength{\headheight}{15pt}
\fancyhf{}
\fancyhead[C]{第一题：二维对流方程完整学习框架}
\fancyfoot[C]{\thepage}
\renewcommand{\headrulewidth}{0.4pt}
\renewcommand{\footrulewidth}{0pt}

\newcolumntype{L}[1]{>{\raggedright\arraybackslash}p{#1}}
\newcommand{\Co}{\mathrm{Co}}
\newcommand{\Cmax}{\mathrm{Co}_{\max}}

\title{\bfseries 第一题：二维对流方程\\完整学习框架与老师的考察目的}
\author{}
\date{}

\begin{document}

\maketitle
\vspace{-2.5em}

\section{总体目标}

第一题不是单纯要求“把求解器跑起来”，而是希望建立一条完整的数值CFD认识链：

\[
\boxed{
\text{方程描述什么}
\longrightarrow
\text{初边值条件确定什么}
}
\]

\[
\boxed{
\text{有限体积如何离散}
\longrightarrow
\text{数值格式引入什么误差}
}
\]

\[
\boxed{
\text{实验参数怎样影响结果}
\longrightarrow
\text{如何根据图像和误差判断原因}
}
\]

最终需要理解的不是“哪个参数得到的误差更小”，而是

\begin{quote}
\centering\bfseries
$N$、网格类型、$\Cmax$、空间格式和时间格式分别改变了离散方程中的什么，\\
为什么最终表现为图上的某种现象。
\end{quote}

\section{原偏微分方程描述什么}

第一题的控制方程为

\begin{equation}
\frac{\partial\phi}{\partial t}
+\nabla\cdot(\bm u\phi)=0.
\end{equation}

其中：

\begin{itemize}
  \item $\phi(x,y,t)$ 是被速度场携带的标量，可以表示染料浓度、污染物浓度、温度标记或体积分数；
  \item $\bm u$ 是给定的速度场；
  \item $\bm u\phi$ 是标量的对流通量。
\end{itemize}

展开散度项：

\begin{equation}
\nabla\cdot(\bm u\phi)
=
\bm u\cdot\nabla\phi
+\phi\nabla\cdot\bm u.
\end{equation}

如果速度场不可压缩，即

\begin{equation}
\nabla\cdot\bm u=0,
\end{equation}

则方程变为

\begin{equation}
\frac{\partial\phi}{\partial t}
+\bm u\cdot\nabla\phi=0.
\end{equation}

沿流体运动轨迹

\begin{equation}
\frac{\mathrm d\bm x}{\mathrm dt}=\bm u
\end{equation}

有

\begin{equation}
\frac{\mathrm d\phi}{\mathrm dt}=0.
\end{equation}

其物理意义是：标量只随流体移动，不会在连续方程中自行变模糊、变矮、变宽或产生振荡。

第一题没有物理扩散项。因此，如果数值结果中出现峰值降低、轮廓展宽、槽口被填充等现象，主要来自数值格式产生的人工误差。

\section{两个算例与方程的关系}

两个算例求解的是同一个对流方程，只是使用了不同的速度场、初始条件、边界条件和终止时间。

\subsection{正弦波平移}

给定

\begin{equation}
\phi(x,y,0)=\sin 2\pi(x+y),
\qquad
\bm u=(1,1).
\end{equation}

精确解为

\begin{equation}
\phi(x,y,t)
=
\sin 2\pi(x+y-2t).
\end{equation}

计算区域为

\begin{equation}
[0,1]^2,
\end{equation}

并在 $x$、$y$ 方向施加周期边界条件。由于正弦波在两个方向均移动一个完整周期，所以

\begin{equation}
\phi(x,y,1)=\phi(x,y,0).
\end{equation}

这个算例主要用于检验：

\begin{itemize}
  \item 光滑解上的数值精度；
  \item 振幅是否衰减；
  \item 相位是否偏移；
  \item $L_1$误差；
  \item 网格收敛阶；
  \item 三角形和四边形网格的差别。
\end{itemize}

\subsection{复杂轮廓刚体旋转}

速度场为

\begin{equation}
\bm u=(0.5-y,\;x-0.5).
\end{equation}

它表示绕

\begin{equation}
(0.5,0.5)
\end{equation}

以单位角速度旋转，所以旋转周期为

\begin{equation}
T=2\pi.
\end{equation}

因此在 $t=2\pi$ 时，所有轮廓理论上都应回到初始位置和初始形状。

\begin{longtable}{L{0.20\linewidth}L{0.32\linewidth}L{0.36\linewidth}}
\toprule
初始结构 & 数学性质 & 主要检验\\
\midrule
\endfirsthead
\toprule
初始结构 & 数学性质 & 主要检验\\
\midrule
\endhead
开槽圆盘 & 存在间断、尖角和窄槽 & 间断分辨率、有界性和数值耗散\\
圆锥 & 连续但存在尖峰和梯度突变 & 峰值保持和轮廓展宽\\
光滑凸峰 & 相对光滑 & 光滑结构上的耗散和相位误差\\
\bottomrule
\end{longtable}

三个形状分别暴露数值格式的不同问题。

\section{初始条件与边界条件}

\subsection{初始条件}

初始条件

\begin{equation}
\phi(x,y,0)=\phi_0(x,y)
\end{equation}

决定被输运的标量一开始是什么形状。同一个求解器可以输运不同初值，不需要为每一个算例重新开发求解器。

\subsection{边界条件}

对于边界外法向量 $\bm n$，考察

\begin{equation}
\bm u\cdot\bm n.
\end{equation}

\begin{itemize}
  \item 若 $\bm u\cdot\bm n<0$，为入流边界，需要规定从外部进入的 $\phi$；
  \item 若 $\bm u\cdot\bm n>0$，为出流边界，通常由内部解决定；
  \item 周期边界把一侧流出的信息送入另一侧。
\end{itemize}

正弦波算例采用

\begin{equation}
\phi(0,y,t)=\phi(1,y,t),
\qquad
\phi(x,0,t)=\phi(x,1,t).
\end{equation}

复杂旋转算例原题没有明确给出边界条件。由于三个轮廓始终远离外边界，可以在入流位置设置背景值

\begin{equation}
\phi=0,
\end{equation}

但需要在报告中明确说明这一处理。

\section{有限体积法如何转化原方程}

对控制体 $\Omega_P$ 积分：

\begin{equation}
\int_{\Omega_P}
\frac{\partial\phi}{\partial t}\,\mathrm dV
+
\int_{\Omega_P}
\nabla\cdot(\bm u\phi)\,\mathrm dV
=0.
\end{equation}

利用高斯散度定理：

\begin{equation}
\int_{\Omega_P}
\nabla\cdot(\bm u\phi)\,\mathrm dV
=
\int_{\partial\Omega_P}
\bm u\phi\cdot\bm n\,\mathrm dS.
\end{equation}

定义单元平均值

\begin{equation}
\phi_P
=
\frac{1}{V_P}
\int_{\Omega_P}\phi\,\mathrm dV,
\end{equation}

将面积分写成各个面的通量之和：

\begin{equation}
V_P\frac{\mathrm d\phi_P}{\mathrm dt}
+
\sum_{f\in\partial\Omega_P}
F_f\phi_f
=0,
\end{equation}

其中

\begin{equation}
F_f=\bm u_f\cdot\bm S_f.
\end{equation}

得到半离散形式：

\begin{equation}
\boxed{
\frac{\mathrm d\phi_P}{\mathrm dt}
=
-\frac{1}{V_P}
\sum_fF_f\phi_f
}.
\end{equation}

这里需要理解：

\begin{itemize}
  \item $V_P$：控制体面积或体积；
  \item $\bm S_f$：带有方向的面面积向量；
  \item $F_f$：通过面 $f$ 的体积通量；
  \item $\phi_f$：面上的标量值；
  \item 单元内标量的变化由所有面通量的净和决定。
\end{itemize}

数学上对应高斯散度定理、积分守恒律、单元平均和方法线，即把PDE转化为有限维ODE系统。

\section{为什么需要构造面值}

有限体积法保存的是单元值

\begin{equation}
\phi_P,\qquad\phi_N,
\end{equation}

但通量通过单元面发生，所以需要确定

\begin{equation}
\phi_f.
\end{equation}

不同空间离散格式的核心区别，就是如何根据周围单元值重构面值。它直接影响空间精度、数值耗散、数值色散、有界性、是否产生负值或过冲，以及对间断结构的分辨能力。

\section{一阶迎风与二阶迎风}

\subsection{一阶迎风}

定义

\begin{equation}
F_f=\bm u_f\cdot\bm S_f.
\end{equation}

一阶迎风面值为

\begin{equation}
\phi_f=
\begin{cases}
\phi_P, & F_f\geq0,\\
\phi_N, & F_f<0.
\end{cases}
\end{equation}

它把单元内的解看成分片常数，因此空间误差通常为

\begin{equation}
E_{\mathrm{space}}=O(h).
\end{equation}

其优点包括稳定、守恒、在合适CFL下通常有界，并且不容易产生负值和强烈振荡；缺点是数值耗散较大。

一维一阶迎风与向前欧拉的修正方程为

\begin{equation}
\phi_t+a\phi_x
=
\frac{a\Delta x}{2}(1-C)\phi_{xx}
+\cdots,
\end{equation}

其中

\begin{equation}
C=\frac{a\Delta t}{\Delta x}.
\end{equation}

右端的二阶导数相当于人工扩散项，因此数值结果会出现：

\begin{itemize}
  \item 正弦波振幅降低；
  \item 开槽圆盘边缘变模糊；
  \item 槽口逐渐被填平；
  \item 圆锥峰值下降、底部变宽；
  \item 光滑凸峰变矮、变宽。
\end{itemize}

\subsection{二阶迎风}

二阶迎风利用上游单元值和梯度重构面值，例如

\begin{equation}
\phi_f
=
\phi_U
+
\nabla\phi_U\cdot
(\bm x_f-\bm x_U).
\end{equation}

在光滑区域通常有

\begin{equation}
E_{\mathrm{space}}=O(h^2).
\end{equation}

与一阶迎风相比，二阶迎风通常能更好地保持正弦波振幅、圆锥峰值和光滑凸峰，并显著减弱数值耗散。

但在间断附近可能产生

\begin{equation}
\phi<0
\qquad\text{或}\qquad
\phi>1,
\end{equation}

以及过冲、欠冲和正负交替波纹。因此通常需要限制器。限制器在光滑区域保留二阶精度，在间断附近减小梯度，必要时退化为一阶迎风。

\section{时间离散为什么会限制空间精度}

向前欧拉格式为

\begin{equation}
\phi_P^{n+1}
=
\phi_P^n
-\frac{\Delta t}{V_P}
\sum_fF_f^n\phi_f^n.
\end{equation}

其时间精度为

\begin{equation}
E_{\mathrm{time}}=O(\Delta t).
\end{equation}

如果空间采用二阶格式，而时间仍使用向前欧拉，并且固定CFL，则

\begin{equation}
\Delta t=O(h).
\end{equation}

总体误差为

\begin{equation}
E
=
O(h^2)+O(\Delta t)
=
O(h^2)+O(h)
=
O(h).
\end{equation}

所以即使空间采用二阶迎风，最终测得的总体收敛阶仍可能接近一阶。

要观察二阶空间精度，可以使用RK2等二阶时间格式并保持固定CFL，或者将向前欧拉时间步缩小到 $\Delta t=O(h^2)$。

\section{\texorpdfstring{$N$}{N}的含义和影响}

首先必须明确程序中的 $N$ 表示每个方向的区间数，还是每条边的顶点数。

如果 $N$ 是每个方向的区间数，则

\begin{equation}
h=\frac1N.
\end{equation}

四边形单元数约为

\begin{equation}
N_{\mathrm{quad}}=N^2,
\end{equation}

三角形单元数约为

\begin{equation}
N_{\mathrm{tri}}=2N^2.
\end{equation}

如果 $N$ 表示顶点数，则每个方向有 $N-1$ 个区间：

\begin{equation}
N_{\mathrm{quad}}=(N-1)^2,
\qquad
N_{\mathrm{tri}}=2(N-1)^2.
\end{equation}

增大 $N$ 会导致

\begin{equation}
N\uparrow
\Longrightarrow
h\downarrow
\Longrightarrow
\text{空间分辨率提高}
\Longrightarrow
\text{耗散减小}
\Longrightarrow
\text{计算成本增加}.
\end{equation}

固定CFL时，

\begin{equation}
\Delta t=O(h)=O(N^{-1}).
\end{equation}

二维单元数约为 $O(N^2)$，时间步数约为 $O(N)$，因此显式求解总工作量大致为

\begin{equation}
O(N^3).
\end{equation}

所以 $N$ 加倍以后，计算量可能增加约八倍。

\section{三角形与四边形网格}

\subsection{四边形网格}

规则四边形网格通常具有面方向规则、单元中心关系简单、单元数量较少和结构化计算方便等特点，但可能表现出明显的网格方向性。

\subsection{三角形网格}

三角形网格通常更适合复杂几何，面方向更多，相同节点间距下单元数量更多，owner-neighbour关系更复杂，并且梯度重构对网格质量更加敏感。

两种网格的面方向、面面积、单元中心和相邻关系不同，而通量取决于

\begin{equation}
F_f=\bm u_f\cdot\bm S_f.
\end{equation}

因此两种网格会产生不同的方向性误差和人工耗散。结果中可能表现为某些方向更加模糊、圆形逐渐变成椭圆形、轮廓出现不对称，或者两种网格收敛阶接近但误差大小不同。

不能预先断定三角形一定比四边形差。比较时应控制近似相同的特征网格尺寸 $h$、格式和CFL。

\section{\texorpdfstring{$\Cmax$}{Co\_max}的含义和影响}

常见局部Courant数为

\begin{equation}
\Co_P
=
\frac{\Delta t}{2V_P}
\sum_f|F_f|,
\end{equation}

最大Courant数为

\begin{equation}
\Cmax
=
\max_P\Co_P.
\end{equation}

它表示一个时间步内，信息相对于当前单元尺度传播了多远。

对规则方形网格和速度

\begin{equation}
\bm u=(1,1),
\end{equation}

若

\begin{equation}
\Delta x=\Delta y=h,
\end{equation}

则大致有

\begin{equation}
\Co
=
\frac{2\Delta t}{h},
\end{equation}

所以

\begin{equation}
\Delta t
\approx
\frac{\Cmax h}{2}.
\end{equation}

题目规定

\begin{equation}
\Cmax=0.2,
\end{equation}

意味着不同网格必须根据各自的 $h$ 调整时间步。

\subsection{\texorpdfstring{$\Cmax$}{Co\_max}较小}

\begin{itemize}
  \item 时间步更小；
  \item 时间步数更多；
  \item 计算成本更高；
  \item 时间离散误差通常减小；
  \item 但一阶迎风的总数值耗散未必减小。
\end{itemize}

\subsection{\texorpdfstring{$\Cmax$}{Co\_max}适当增大}

在稳定范围内可能减少计算步数、减弱一阶迎风的部分数值耗散，但会增大时间截断误差并降低稳定裕度。

\subsection{\texorpdfstring{$\Cmax$}{Co\_max}超过稳定限制}

可能出现数值振荡、负值和过冲、棋盘状噪声、最大值快速增长并最终发散。

\section{结果图中的现象与原因}

\begin{longtable}{L{0.39\linewidth}L{0.49\linewidth}}
\toprule
结果现象 & 主要原因\\
\midrule
\endfirsthead
\toprule
结果现象 & 主要原因\\
\midrule
\endhead
正弦波振幅降低 & 数值耗散\\
波峰位置偏移 & 数值色散或时间相位误差\\
开槽圆盘边缘模糊 & 一阶迎风人工扩散\\
槽口被填平 & 高值向槽内人工扩散\\
圆锥峰值下降、底部变宽 & 人工扩散重新分配质量\\
光滑凸峰变矮、变宽 & 数值耗散\\
间断附近出现波纹 & 无限制高阶重构产生色散\\
出现负值或超过最大值 & 格式无界、限制器不足或CFL过大\\
轮廓旋转角度不正确 & 时间积分、速度场或终止时间错误\\
圆形变成椭圆形 & 网格各向异性耗散\\
旋转半径改变 & 速度插值或面通量计算错误\\
总质量持续改变 & 面通量不守恒或边界泄漏\\
边界附近出现异常 & 周期边界或入流边界处理错误\\
网格加密后误差不下降 & 时间误差、边界误差或实现错误主导\\
CFL增大后突然出现噪声 & 接近或超过显式稳定上限\\
\bottomrule
\end{longtable}

\section{后处理需要观察什么}

\subsection{正弦波算例}

需要观察：

\begin{enumerate}
  \item 数值解云图；
  \item 精确解云图；
  \item 有符号误差图；
  \item 绝对误差图；
  \item 数值解与精确解的一维截线；
  \item $L_1$误差；
  \item 不同 $N$ 下的收敛阶；
  \item 最大值和最小值；
  \item 三角形与四边形网格的结果。
\end{enumerate}

其中，振幅降低对应耗散，波峰横向偏移对应色散，误差随 $N$ 的变化反映收敛性。

\subsection{复杂轮廓算例}

需要观察：

\begin{enumerate}
  \item 初始等值线；
  \item $t=2\pi$时的最终等值线；
  \item 初始与最终轮廓叠加；
  \item 差值图；
  \item 总质量随时间的变化；
  \item 最大值和最小值随时间的变化；
  \item 三个轮廓的一维截线；
  \item 三角形与四边形网格对比；
  \item 一阶与二阶迎风对比。
\end{enumerate}

所有对比图必须使用相同的色标范围、等值线级别、坐标范围、纵横比例和输出时刻。

\section{真正需要完成的四类对比}

\subsection{\texorpdfstring{$N$}{N}的影响}

固定网格类型、空间格式、时间格式和

\begin{equation}
\Cmax=0.2,
\end{equation}

只改变 $N$。需要理解

\begin{equation}
N\uparrow
\Longrightarrow
h\downarrow
\Longrightarrow
\text{分辨率提高、耗散减小、成本增加}.
\end{equation}

\subsection{网格类型的影响}

固定近似相同的 $h$、格式、CFL和初边值条件，比较三角形与四边形。需要理解

\begin{equation}
\text{面方向与连接关系不同}
\Longrightarrow
\text{数值通量不同}
\Longrightarrow
\text{方向性耗散不同}.
\end{equation}

\subsection{一阶与二阶迎风的影响}

固定其他条件，比较

\begin{equation}
\text{一阶迎风}
\Longrightarrow
\text{耗散较大、稳定、有界},
\end{equation}

以及

\begin{equation}
\text{二阶迎风}
\Longrightarrow
\text{光滑区更准确、耗散较小、间断附近可能振荡}.
\end{equation}

同时判断一阶时间格式是否限制了二阶空间格式的总体收敛阶。

\subsection{\texorpdfstring{$\Cmax$}{Co\_max}的影响}

固定网格和格式，只改变CFL。需要理解

\begin{equation}
\Cmax
\Longrightarrow
\Delta t
\Longrightarrow
\text{稳定性、时间误差、数值耗散与计算成本}.
\end{equation}

\section{建议学习顺序}

\begin{enumerate}
  \item 理解一维对流方程和特征线；
  \item 理解连续解为什么只平移、不扩散；
  \item 推导一维有限体积半离散形式；
  \item 推导一阶迎风和向前欧拉；
  \item 推导CFL稳定条件；
  \item 用修正方程理解数值耗散；
  \item 编写一维正弦波求解器；
  \item 比较不同 $N$ 和不同CFL；
  \item 扩展到二维四边形网格；
  \item 完成二维正弦波误差和收敛阶；
  \item 学习面、owner、neighbour和面积向量；
  \item 扩展到三角形非结构网格；
  \item 学习二阶梯度重构和限制器；
  \item 比较一阶与二阶迎风；
  \item 完成复杂轮廓刚体旋转；
  \item 从结果图判断耗散、色散、网格方向性、守恒和稳定性。
\end{enumerate}

\section{老师发这道题的目的}

从题目的结构看，老师的目的确实包括让学生理解上述内容，而不只是完成代码。

\begin{itemize}
  \item 要求写半离散形式，说明要理解方程到有限体积通量守恒的转化；
  \item 要求开发显式求解器，说明要理解空间面值、时间推进和CFL；
  \item 要求逐渐加密网格并计算误差和收敛阶，说明要理解 $N$、$h$、$\Delta t$和精度的关系；
  \item 要求分别使用三角形和四边形网格，说明要理解非结构网格和网格拓扑的影响；
  \item 要求计算复杂轮廓，说明要观察数值耗散、有界性和间断分辨能力；
  \item 指定CFL，说明时间步不是任意参数，而与网格和稳定性相关。
\end{itemize}

老师最低限度可能希望看到：

\begin{itemize}
  \item 正确推导半离散方程；
  \item 独立构建可运行求解器；
  \item 完成三角形和四边形计算；
  \item 得到合理的结果图；
  \item 误差随网格加密而下降；
  \item 能解释最主要的数值耗散现象。
\end{itemize}

如果进一步能够解释 $N$为什么影响耗散和计算成本、$\Cmax$为什么影响稳定性和误差、一阶与二阶迎风为什么表现不同、振幅误差和相位误差如何区分、三角形与四边形为什么存在方向性差异，以及为什么一阶时间格式可能限制二阶空间格式，就说明已经不只是完成了练习，而是在建立开发和验证CFD求解器所需的数值判断能力。

\end{document}
